\chapter*{Preface}
\addcontentsline{toc}{chapter}{Preface}
\markboth{Preface}{}

About three years ago the questions from our customers changed.

I am a
product manager at Braiins, working on Braiins Manager, the software that runs
mining fleets, and for years their requests had only been about hashrate. Then they
became about when to switch it off. Our customers wanted their machines to slow down when
prices spiked, and to know whether they should be selling reserves to the grid
operator. I knew what a joule per terahash was. I had no
idea what an ancillary service was, or why the same megawatt-hour could cost
twenty-one dollars before dawn and nearly four thousand that afternoon. % src: analysis/data/ercot_rtm15_spp_loadzones_2021on.csv.gz (ERCOT MIS 13061, Central prevailing time), LZ\_WEST 15-minute real-time settlement point prices on November 17, 2024: \$21.46 in the 04:30 interval, the lowest before the spike, and \$3,925.09 in the 15:00 interval, the day's high (the day's low, \$13.29, came at 23:00); ercot2026protocols sec. 6.6.1.2 and 6.6.3.2 (a load zone's energy settles at a 15-minute price built from the five-minute LMPs: MIS 13061 publishes the Real-Time Settlement Point Price of 6.6.1.2(1), and a load's own charge in 6.6.3.2(2) uses its energy-weighted form, 6.6.1.2(2); the formula is the one in force since RTC+B, and the 2024 price also carried the ORDC adder that RTC+B removed, Chapter 3). The day's five-minute LMPs (analysis/data/ercot_spp_rtm5_lzwest.csv.gz, NP6-788-CD, stamped in UTC) ran from \$7.00 at 00:10 to \$4,979.08 at 15:10 Central; this line printed those five-minute values as what a megawatt-hour cost until 2026-09-26, and read them as morning and evening until 2026-09-25

Answering those questions meant learning a second industry. I found two professional worlds describing the same physical event in completely different words: A grid operator sees a load responding to a price
signal. A miner sees a site pausing to protect its margin.

So I went looking for the book that would translate between them. For the
Bitcoin half of the problem there is \textit{Mastering Bitcoin}, Andreas
Antonopoulos's book that takes a working engineer and makes the protocol
legible. For the power-system half there are textbooks: excellent, forbidding,
and silent on mining. In between, where my actual question lived, I found
conference slides, vendor decks, and a great many confident numbers with
nothing behind them. I could not find that book.

Working it out myself took longer than it should have. I read market protocols
written for traders, pulled the settlement files ERCOT (the Electric Reliability Council of Texas) publishes, and asked operators questions that must have sounded naive. The
connections came slowly, and usually in the wrong order. This book is the
write-up of that work, arranged the way I wish somebody had handed it to me.

I only trust something once I can compute it. That is as much a limitation as a
method, and it shaped every chapter. Rather than repeat an industry figure for
what flexibility is worth, I pulled the price history and worked it out hour by
hour, in each ERCOT load zone. Where a number in these pages
is mine, the script that produced it ships with the book, in a public repository
at \repourl. Where it belongs to
somebody else, the citation says whose. If a result bothers you, you do not
have to argue with me about it. Open the script, change the assumption you
object to, and run it again.

You should also know where I sit. Braiins sells Bitcoin mining firmware and fleet
software and runs a \glossaryterm{mining-pool}{mining pool}, a service that combines miners' computing power and shares the proceeds, so I have an obvious interest in mining being
taken seriously as a grid resource. The measured machine efficiency figures in these
chapters come from aggregated operating data from miners running Braiins~OS,
because that is real data I can inspect at scale rather than a manufacturer's rated number. I say so every
time I use it. My answer to that conflict of interest is to publish the arithmetic and let you check it.

The book has two limits. The first is the territory. I chose Texas because that is
where flexible mining load has gone furthest, and for a more selfish reason:
ERCOT publishes enough data that I could check my own work. The mechanisms
travel further than the numbers do, and the book is built so that other markets
can be added later on the same foundation. The second limit is me. I am a
product manager, not a power system engineer and not an economist. Where the
data carries a claim, I say so. Where I am making a judgment call, I mark it as
one.

I am finishing these pages in a hospital room, where my daughter was born this
week and where I am staying with her and her mother.

\vspace{1em}
\noindent Vit Ficl \\
September 9, 2026
